\section{Open challenges: an evidence-chained gap register}
\label{sec:challenges}

Each gap below is grounded in at least two audited references, annotated with the database that
surfaced the evidence (arXiv / OpenAlex / Crossref, per the project search log), and closed
with the reasoning why current work leaves it open. The register is maintained as a standalone
deliverable (\texttt{notes/research\_gaps.md}) and summarized here.

\textbf{G1 --- No standardized cross-family HER benchmark.} Factor deconvolution exists within
a single family (S-COF/FS-COF; OpenAlex)~\cite{ghosh2020identification}, while performance is
reported in incommensurable currencies --- AQE at a stated wavelength
(OpenAlex)~\cite{li2022covalent} versus mass-normalized rates
(OpenAlex)~\cite{yang2021protonated} --- and high-throughput screening has targeted H$_2$O$_2$
rather than HER (OpenAlex)~\cite{zhao2022accelerated}. \emph{Why open:} multi-variable control
samples are costly, and reporting conventions differ across groups --- the same library
contains AQE-denominated and mass-rate-denominated headline results with no conversion
basis~\cite{li2022covalent,yang2021protonated}; no library study runs two linkage families
under one fixed protocol. \emph{Needed:} a fixed-protocol,
multi-family benchmark reporting both AQE and mass rate.

\textbf{G2 --- Overall water splitting remains a single-example regime.} The only library
system targeting non-sacrificial operation is the chemically bonded 2D/2D COF/WO$_3$ Z-scheme,
with rates far below sacrificial benchmarks (OpenAlex)~\cite{shen2023in}; the 2020 CSR review
confirms sacrificial dominance across COF photocatalysis (OpenAlex/Crossref)~\cite{wang2020covalent};
even molecular-cocatalyst systems rely on sacrificial donors (OpenAlex)~\cite{banerjee2017single}.
\emph{Why open:} sacrificial donors mask the four-electron oxidation half-reaction, whose
oxidative stress and band-position demands most COFs were never designed to meet.
\emph{Needed:} valence-band and anti-oxidation engineering plus Z-scheme charge routing
designed for the OER half-reaction.

\textbf{G3 --- Operando structural durability is unquantified.} Coaxial stacking disorders
during photocatalytic cycling and activity decays unless the stack is locked
(OpenAlex)~\cite{zhou2021peg}; chemical stability tests (9~N HCl, boiling water;
Crossref)~\cite{kandambeth2012construction} probe a different failure axis, and the
linkage-design trilemma contains no operando dimension (OpenAlex)~\cite{haase2020solving}.
\emph{Why open:} the field equates chemical with photocatalytic stability, and standard
practice is before/after PXRD rather than time-resolved structural tracking. \emph{Needed:}
operando structure--activity correlation (cycles vs.\ crystallinity vs.\ rate) and a shared
``photocatalytic lifetime'' reporting convention.

\textbf{G4 --- Green and scalable routes are not HER-benchmarked.} Flow-derived TpPa-1 reports
BET area, space-time yield, and energy metrics but no HER data
(Crossref)~\cite{xu2026structural}; the microwave comparison stops at porosity and CO$_2$
uptake (OpenAlex)~\cite{grenu2020microwave,wei2015the}; room-temperature and flow routes close
with crystallinity and porosity (OpenAlex)~\cite{peng2016room}. \emph{Why open:} the synthesis
methodology and photocatalysis communities publish to different completion criteria, so the
route$\rightarrow$crystallinity$\rightarrow$HER chain --- whose middle link is causally
established (OpenAlex)~\cite{ghosh2020identification} --- has never been closed on one
material. \emph{Needed:} same-protocol HER comparison of solvothermal, microwave, flow, and
mechanochemical TpPa-1 --- precisely the experiment Route~C (Sec.~\ref{sec:routec}) is designed
to seed.

\textbf{G5 --- Exciton-binding-energy design rules lack an experimental loop.} DFT screening
makes $E_b$ tunable on paper (OpenAlex)~\cite{qian2023computation}; excited-state mapping is
family-specific (OpenAlex)~\cite{chen2023tuning}; and the mechanistic perspective explicitly
calls the pipeline's characterization methodology unsettled (OpenAlex)~\cite{blatte2024photons}.
\emph{Why open:} measuring $E_b$ on polycrystalline powders is noisy and mismatched to idealized
computational models, and computational and spectroscopic groups rarely share sample sets.
\emph{Needed:} a computed-then-measured $E_b$ benchmark across one isoreticular series.

\textbf{G6 --- Earth-abundant cocatalysts trail Pt by orders of magnitude.} The
cobaloxime/N$_2$-COF system achieves 782~$\mu$mol\,h$^{-1}$\,g$^{-1}$ at TON 54.4
(OpenAlex)~\cite{banerjee2017single}; covalent wiring prolongs but does not eliminate
degradation (OpenAlex)~\cite{gottschling2020rational}; meanwhile the field's effort
concentrates on Pt speciation (OpenAlex)~\cite{dong2021platinum,li2022in}. \emph{Why open:}
molecular-cocatalyst failure modes demand operando spectro-electrochemistry beyond most COF
groups' stacks, while Pt photodeposition is operationally trivial. \emph{Needed:} self-healing
ligand design on covalent anchors, and lifecycle TON accounting that makes the Pt/non-Pt
comparison honest.

These six gaps are not independent: G1's protocol deficit blocks progress measurement on G2--G6,
and G4 is the gap this survey's Route~C plan directly operationalizes. The gap register with
full evidence chains and database provenance accompanies this survey as an audited deliverable.
